% =====================================================================
%  Midterm Thesis Report: Improving Regional Subseasonal-to-Seasonal
%  (S2S) Forecasting.  Overleaf-ready: upload this folder as-is
%  (compiler: pdfLaTeX).  Figures are rebuilt by the scripts in tools/.
% =====================================================================
\documentclass[11pt,a4paper]{article}

% ---------------- Cover-page fields (edit these) ---------------------
\newcommand{\StudentName}{Ayush Raj}
\newcommand{\RegNo}{20221068}
\newcommand{\Programme}{BS-MS Dual Degree Programme}
\newcommand{\Department}{Department of Data Science}
\newcommand{\SupName}{Prof.\ Sandeep Kumar Juneja}
\newcommand{\SupAffilA}{Director, Safexpress Centre for Data,}
\newcommand{\SupAffilB}{Learning and Decision Sciences}
\newcommand{\SupAffilC}{Ashoka University}
\newcommand{\ExpName}{Dr.\ Bedartha Goswami}  % CONFIRM: expert at IISER Pune
\newcommand{\ExpAffilA}{Assistant Professor}
\newcommand{\ExpAffilB}{Department of Data Science}
\newcommand{\ExpAffilC}{IISER Pune}
\newcommand{\ReportDate}{October 2026}
% ---------------------------------------------------------------------

\usepackage[T1]{fontenc}
\usepackage[utf8]{inputenc}
\usepackage{amsmath}
\usepackage{newpxtext,newpxmath}
\usepackage{microtype}
\usepackage{soul}
\sodef\spaced{}{0.14em}{0.5em}{0em}   % letter-spaced small caps (works in pdfLaTeX and XeLaTeX)
\usepackage[margin=1.9cm,top=1.9cm,bottom=1.9cm,headsep=0.45cm]{geometry}
\usepackage{graphicx}
\usepackage{booktabs,tabularx,multirow,multicol,array}
\usepackage[table]{xcolor}
\usepackage{enumitem}
\usepackage{titlesec}
\usepackage{needspace}
\usepackage{caption}
\usepackage{fancyhdr}
\usepackage[numbers,sort&compress]{natbib}
\usepackage[colorlinks=true,linkcolor=black,citecolor=black,urlcolor=black]{hyperref}
\urlstyle{same}

\definecolor{accent}{HTML}{0F3A63}
\definecolor{accentmid}{HTML}{2A6496}
\definecolor{accentlight}{HTML}{EAF1F8}
\definecolor{warm}{HTML}{C0392B}
\definecolor{warmlight}{HTML}{FBEDEB}
\definecolor{ok}{HTML}{1E7D4F}
\definecolor{inkgrey}{HTML}{55606B}
\newcommand{\placeholder}[1]{\colorbox{yellow!45}{\textit{[#1]}}}
\definecolor{ink}{HTML}{1B2A41}  % single muted accent for the title page
\newcommand{\ornament}{{\color{ink}\rule[0.5ex]{1.6cm}{0.4pt}\hspace{0.35cm}\raisebox{0.15ex}{\scriptsize$\blacklozenge$}\hspace{0.35cm}\rule[0.5ex]{1.6cm}{0.4pt}}\par}

\setlength{\parskip}{2pt}
\setlength{\textfloatsep}{10pt plus 2pt minus 2pt}
\setlength{\intextsep}{8pt plus 2pt minus 2pt}
\setlength{\parindent}{1.2em}
\setlist{nosep,leftmargin=1.6em}
\captionsetup{font=small,labelfont=bf,skip=4pt}
\captionsetup[table]{position=top,skip=3pt}
\titleformat{\section}{\normalfont\large\bfseries}{\thesection}{0.6em}{}
\titleformat{\subsection}{\normalfont\normalsize\bfseries}{\thesubsection}{0.6em}{}
\titlespacing*{\section}{0pt}{9pt}{3pt}
\titlespacing*{\subsection}{0pt}{5pt}{1pt}

\pagestyle{fancy}
\fancyhf{}
\renewcommand{\headrulewidth}{0.4pt}
\fancyhead[L]{\footnotesize Mid-Year Thesis Report}
\fancyhead[R]{\footnotesize \StudentName}
\fancyfoot[C]{\small\thepage}


\begin{document}

% ============================ COVER =================================
\begin{titlepage}
\thispagestyle{empty}
\newgeometry{margin=2.6cm,top=2.3cm,bottom=2.3cm}
\centering
% --- institutional marks
\begin{minipage}[c]{0.3\linewidth}\raggedright
  \includegraphics[height=2.0cm]{logos/iiserpune.png}
\end{minipage}\hfill
\begin{minipage}[c]{0.66\linewidth}\raggedleft
  \includegraphics[height=1.9cm]{logos/ashoka.png}\hspace{0.18cm}%
  \raisebox{0.28cm}{\includegraphics[height=1.32cm]{logos/scdlds.png}}
\end{minipage}

\vspace{0.35cm}
{\color{ink}\rule{\linewidth}{0.4pt}}

\vspace{3.2cm}
{\small\color{ink}\scshape\spaced{Mid-Year Thesis Report}\par}
\vspace{0.6cm}
{\fontsize{26}{33}\selectfont\color{ink}
Improving Regional\\[2pt] Subseasonal-to-Seasonal Forecasting\par}
\vspace{0.55cm}
\ornament

\vspace{2.4cm}
{\Large\scshape \StudentName\par}
\vspace{0.2cm}
{\small Registration No.\ \RegNo\par}
{\small \Department, IISER Pune\par}
{\small Safexpress Centre for Data, Learning and Decision Sciences, Ashoka University\par}

\vspace{2.4cm}
\renewcommand{\arraystretch}{1.2}
\begin{tabular}{@{}>{\centering\arraybackslash}p{0.45\linewidth}!{\color{black!35}\vrule width 0.4pt}>{\centering\arraybackslash}p{0.45\linewidth}@{}}
{\footnotesize\color{ink}\scshape\spaced{Supervisor}} & {\footnotesize\color{ink}\scshape\spaced{Expert}} \\[1pt]
\SupName & \ExpName \\
{\footnotesize\itshape \SupAffilA} & {\footnotesize\itshape \ExpAffilA} \\
{\footnotesize\itshape \SupAffilB} & {\footnotesize\itshape \ExpAffilB} \\
{\footnotesize\itshape \SupAffilC} & {\footnotesize\itshape \ExpAffilC} \\
\end{tabular}



\vfill
{\color{ink}\rule{\linewidth}{0.4pt}}

\vspace{0.3cm}
{\normalsize Indian Institute of Science Education and Research Pune\par}
\vspace{0.12cm}
{\small\color{ink}\scshape\expandafter\spaced\expandafter{\ReportDate}\par}
\restoregeometry
\end{titlepage}

% ============================ ABSTRACT ===============================
\setcounter{page}{1}
\begin{center}
{\large\bfseries Abstract}
\end{center}
\vspace{-2pt}
{\small\sloppy\noindent
Subseasonal-to-seasonal (S2S) forecasts, two to six weeks ahead, are needed for agriculture, water management and disaster
preparedness in India, but their skill over the region is not well known and their ensembles are often unreliable. This thesis
studies how such forecasts can be evaluated and improved for India. In the first half of the project we built
IndiaS2S-Bench, a benchmark of nine operational, AI and hybrid forecast systems over India; developed a lightweight regional adapter
that calibrates the FuXi-S2S ensemble and lowers its CRPS by 16.2\% for monsoon rainfall; tested learned multi-model combination;
and set up a real-time experimental forecast service at \url{https://s2s.scdlds.com}. Two papers have been accepted at a NeurIPS 2026
workshop and a third is in preparation. The remaining work will develop a generative deep-learning model for joint probabilistic
forecasts of rainfall and temperature.\par}

\vspace{4pt}
\noindent{\color{black!35}\rule{\linewidth}{0.4pt}}

% ============================ 1. INTRODUCTION ========================
\section{Introduction}
Subseasonal-to-seasonal (S2S) forecasts cover the range from about two weeks to two months ahead~\cite{vitart2017s2s}. This
range lies between weather forecasts and seasonal outlooks. It matters for India because most of the annual rainfall falls during
the southwest monsoon (June--September, JJAS), and a dry or wet spell of a few weeks can affect sowing, reservoir operation and
flood preparation~\cite{imd2022lpa}.

Several AI weather models have been released in recent years, for example GraphCast, FuXi-S2S and
NeuralGCM~\cite{lam2023graphcast,chen2024fuxis2s,kochkov2024neuralgcm}. Most are evaluated on global averages~\cite{rasp2024wb2}.
A model that does well globally need not do well over India, and its skill can change with region and season~\cite{malik2024s2s}.
Rainfall is the most difficult variable at these lead times, so we work with weekly means and anomalies from climatology, and we use
climatology as the reference forecast.

\paragraph{Objectives.} The thesis addresses three questions:
\begin{enumerate}
\item Which forecast system is most skilful over India for a given variable, lead time, region and season?
\item Can a small regional model correct the errors of a global AI forecast, in particular its ensemble spread?
\item Can deep generative models give reliable joint probabilistic forecasts of rainfall and temperature?
\end{enumerate}
Questions~1 and~2 are largely answered (Section~\ref{sec:progress}); Question~3 is the main task for the rest of the
project (Section~\ref{sec:tasks}).

% ============================ 2. DATA AND METHODS ===================
\section{Data and methods}

\paragraph{Forecast systems and observations.} Table~\ref{tab:data} lists the data. Forecasts from four operational centres were
taken from the S2S database; five AI and hybrid models were run by us from ERA5 initial conditions. All models were scored only on
dates after their training period. We used 517 start dates from 2020 to 2024 (516 for temperature), and aggregated forecast days
1--42 into six weekly means on a common 1.5$^\circ$ grid.

\begin{table}[htb]
\centering\small
\caption{Forecast systems and verification data used in this work.}
\label{tab:data}
\begin{tabularx}{\linewidth}{@{}l >{\raggedright\arraybackslash}X >{\raggedright\arraybackslash}p{0.36\linewidth}@{}}
\toprule
Category & Products & Notes \\
\midrule
Operational models & ECMWF, NCEP, UKMO, CMA & S2S database~\cite{vitart2017s2s} \\
AI and hybrid models & FuXi-S2S, NeuralGCM, FourCastNet~3, DLESyM-v0, DLESyM-v1 & Run by us from ERA5 initial conditions \\
Rainfall truth & IMERG; IMD gridded rainfall & Satellite~\cite{huffman2023imerg}; gauge analysis~\cite{pai2014imd} \\
Temperature truth & ERA5 2\,m temperature & Reanalysis~\cite{hersbach2020era5} \\
Regions & All-India and four IMD homogeneous regions & Northwest, Central, South Peninsula, East \& Northeast \\
\bottomrule
\end{tabularx}
\end{table}

\paragraph{Verification.} Deterministic skill is measured with the anomaly correlation coefficient (ACC) and the root-mean-square error
(RMSE) of the ensemble mean. Probabilistic skill is measured with the continuous ranked probability score (CRPS) and the coverage of the
central 90\% ensemble range, which should contain 90\% of observations in a reliable ensemble. Scores are reported for All-India, the four
regions and four seasons: January--February (JF), March--May (MAM), JJAS and October--December (OND).

% ============================ 3. PROGRESS ===========================
\section{Progress to date}\label{sec:progress}

\subsection{IndiaS2S-Bench: comparing forecast systems over India}
Table~\ref{tab:acc} gives the All-India ACC by lead week for four representative systems~\cite{raj2026bench}. The six-model
mean (MME) has the highest Week-1 ACC for rainfall (0.622) and temperature (0.596), and the lowest RMSE in 11 of 12 target--lead cases.
Rainfall skill falls quickly: after Week~2 it is low, and the MME is no longer better than climatology in RMSE.

\begin{table}[htb]
\centering\small
\caption{All-India anomaly correlation by lead week, 2020--2024. Rainfall is verified against IMERG and 2\,m temperature against ERA5.
Bold marks the best of the four systems shown; ECMWF Week-3 rainfall is missing from the archive.}
\label{tab:acc}
\setlength{\tabcolsep}{5.2pt}
\begin{tabular}{@{}l cccccc c cccccc@{}}
\toprule
 & \multicolumn{6}{c}{Rainfall} & & \multicolumn{6}{c}{2\,m temperature} \\
\cmidrule(lr){2-7}\cmidrule(lr){9-14}
Week & 1 & 2 & 3 & 4 & 5 & 6 & & 1 & 2 & 3 & 4 & 5 & 6 \\
\midrule
MME      & \textbf{0.622} & \textbf{0.399} & \textbf{0.228} & 0.136 & 0.098 & 0.057 & & \textbf{0.596} & \textbf{0.471} & \textbf{0.329} & \textbf{0.224} & \textbf{0.168} & \textbf{0.142} \\
FuXi-S2S & 0.586 & 0.354 & 0.213 & 0.151 & \textbf{0.153} & \textbf{0.109} & & 0.559 & 0.409 & 0.269 & 0.190 & 0.139 & 0.123 \\
ECMWF    & 0.569 & 0.382 & --    & \textbf{0.175} & 0.141 & 0.096 & & 0.512 & 0.393 & 0.311 & 0.223 & \textbf{0.168} & \textbf{0.142} \\
UKMO     & 0.553 & 0.313 & 0.161 & 0.054 & 0.051 & 0.045 & & 0.437 & 0.310 & 0.177 & 0.108 & 0.080 & 0.072 \\
\bottomrule
\end{tabular}
\end{table}

The ranking changes with region and season (Figure~\ref{fig:selection}). For JJAS rainfall at Weeks~3--6, FuXi-S2S has the highest ACC in
19 of 20 region--lead cases against IMERG and in 18 of 20 against IMD; for temperature at long leads, ECMWF is the best single model.
Different scores also rank the models differently, so no single system is best for every use. Against IMD's operational ERPAS system on
101 matched cases from 2023--2024, FuXi-S2S has higher ACC in Weeks~1--4, while ERPAS is better in October--December.

\begin{figure}[htb]
\centering
\includegraphics[width=\linewidth]{figures/map_selection.pdf}
\caption{System with the highest ACC in each IMD homogeneous region, by season (rows) and lead week (columns), 2020--2024.
Hatching marks regions where even the best system has ACC below 0.10~\cite{raj2026bench}.}
\label{fig:selection}
\end{figure}

\subsection{Calibrating FuXi-S2S with a regional adapter}
Since FuXi-S2S is the strongest model for long-lead monsoon rainfall, we tried to improve it rather than build a new model. Its 50-member
ensemble is too narrow over India: in Week~1 only 15.6\% of observations fall inside its 90\% range (Table~\ref{tab:adapter}b). We trained
a small neural network with 42,434 parameters, the regional adapter~\cite{raj2026adapter}. It takes the FuXi-S2S members as input and
predicts a shift $\delta$ and a spread factor $s$ for every grid point and lead. These are applied to each member $m$ after a log
transform, $u_m=\log(1+x_m)$:
\begin{equation}
q_m = u_m + \delta + (s-1)\,(u_m-\bar u),
\end{equation}
where $\bar u$ is the ensemble mean. The adapter was trained on FuXi-S2S reforecasts for 2002--2017 against IMD rainfall by minimising
the CRPS, with 2018--2019 used for model selection.

\begin{table}[htb]
\centering\small
\caption{Regional adapter results for JJAS rainfall~\cite{raj2026adapter}. (a) Scores on the 2022--2024 monsoons (600 case--leads);
moment calibration is a simple statistical baseline. (b) Fraction of observations inside the central 90\% ensemble range
(nominal 0.90), 2020--2024, 169 cases per lead.}
\label{tab:adapter}
\begin{minipage}[t]{0.53\linewidth}\centering
\textit{(a) Skill scores}\par\vspace{3pt}
\begin{tabular}{@{}lccc@{}}
\toprule
 & Raw & Moment & Adapter \\
\midrule
CRPS vs IMD (mm\,day$^{-1}$)   & 3.157 & 2.763 & \textbf{2.645} \\
CRPS vs IMERG (mm\,day$^{-1}$) & 3.152 & 2.826 & \textbf{2.767} \\
ACC vs IMD                     & 0.296 & --    & \textbf{0.322} \\
\bottomrule
\end{tabular}
\end{minipage}\hfill
\begin{minipage}[t]{0.44\linewidth}\centering
\textit{(b) 90\% coverage}\par\vspace{3pt}
\begin{tabular}{@{}lccc@{}}
\toprule
Week & Raw & Moment & Adapter \\
\midrule
1 & 0.156 & 0.566 & \textbf{0.686} \\
2 & 0.331 & 0.690 & \textbf{0.761} \\
4 & 0.657 & 0.785 & \textbf{0.816} \\
6 & 0.741 & 0.821 & \textbf{0.841} \\
\bottomrule
\end{tabular}
\end{minipage}
\end{table}

The adapter lowers the CRPS by 16.2\% against IMD and by 12.2\% against IMERG, and the ACC also improves slightly
(Table~\ref{tab:adapter}a). Coverage rises at every lead, most strongly in Weeks~1--2, but stays below the nominal 90\%. The gains are
spread over almost all of India and are largest over the Western Ghats and the Northeast, where rainfall is heaviest
(Figure~\ref{fig:adapter}). A separate U-Net model for the ensemble mean reduced the RMSE by 8.5\% and raised the ACC from 0.276 to
0.364. So far the method has been tested for one model and one variable.

\begin{figure}[htb]
\centering
\includegraphics[width=\linewidth]{figures/map_adapter.pdf}
\caption{Reduction in CRPS from the regional adapter for JJAS rainfall, 2020--2024 (169 cases per lead); positive values mean the
calibrated ensemble is better. Shading interpolates the native 1.5$^\circ$ cells for display~\cite{raj2026adapter}.}
\label{fig:adapter}
\end{figure}

\subsection{Learned combination of models}
We also tested whether a learned combination of seven models can beat their simple average. We used PiggyCast, a stacking method based on
XGBoost~\cite{kimani2025piggycast,chen2016xgboost}, trained on 2020--2023, tuned on 2024 and tested once on the 2025 monsoon against IMD
rainfall~\cite{raj2026piggy}. PiggyCast improves both ACC and RMSE over the equal-weight mean (Table~\ref{tab:piggy}), but an additive bias
correction of the equal-weight mean, tried afterwards, gives a similar result, and the gain is limited to Weeks~1--4. With only a few years
of training data, simple methods are hard to beat, and every new method must be compared with them.

\begin{table}[htb]
\centering\small
\caption{Multi-model combination for JJAS 2025 rainfall against IMD (35 start dates, mean over six leads)~\cite{raj2026piggy}.
Errors are in mm\,day$^{-1}$; $^{\dagger}$exploratory baseline added after the main test.}
\label{tab:piggy}
\begin{tabular}{@{}lcccc@{}}
\toprule
Method & ACC $\uparrow$ & RMSE $\downarrow$ & MAE $\downarrow$ & Bias \\
\midrule
Equal-weight mean                  & 0.307 & 5.174 & 3.715 & $-0.428$ \\
PiggyCast                          & 0.347 & \textbf{4.989} & \textbf{3.603} & $-0.303$ \\
Bias-corrected mean$^{\dagger}$    & \textbf{0.369} & 5.017 & 3.635 & $-0.325$ \\
\bottomrule
\end{tabular}
\end{table}

\subsection{Real-time forecast service}
We have set up an experimental forecast website, \url{https://s2s.scdlds.com}, hosted at SCDLDS. It shows weekly rainfall and 2\,m
temperature forecasts, anomalies and tercile probabilities for Weeks~1--6 from a 100-member FuXi-S2S ensemble, and is updated twice a
week. Anomalies are taken relative to the model's own 2002--2021 reforecast climatology, so that systematic model bias does not appear
as a forecast signal (Figure~\ref{fig:forecast}). The site is for research and testing, not an official forecast.

\begin{figure}[htb]
\centering
\includegraphics[width=\linewidth]{figures/map_forecast.pdf}
\caption{Real-time forecast issued on 27 September 2026 at \url{s2s.scdlds.com}: ensemble-mean weekly rainfall (top) and 2\,m
temperature (bottom) anomalies for Weeks~1--4, from 100 FuXi-S2S members started from ECMWF operational analyses.}
\label{fig:forecast}
\end{figure}

\subsection{Publications}
\begin{enumerate}[label={[P\arabic*]},leftmargin=2.4em]
\item IndiaS2S-Bench: a comprehensive benchmark of subseasonal to seasonal forecasting over India. \textit{Tackling Climate Change with
Machine Learning workshop, NeurIPS 2026} (accepted)~\cite{raj2026bench}.
\item Calibrating AI subseasonal to seasonal forecasts with lightweight regional adapters. \textit{Tackling Climate Change with Machine
Learning workshop, NeurIPS 2026} (accepted)~\cite{raj2026adapter}.
\item When does lightweight post-processing help? A seven-system JJAS rainfall benchmark over India. In preparation~\cite{raj2026piggy}.
\end{enumerate}

% ============================ 4. REMAINING WORK =====================
\section{Remaining work}\label{sec:tasks}
The adapter can only shift and scale the existing ensemble, and the statistical methods above do not model the joint behaviour of
rainfall and temperature. The second half of the project therefore moves to deep generative post-processing. We will train a conditional
generative model based on flow matching~\cite{lipman2023flow,liu2023rectified} that produces rainfall and temperature fields together,
given the FuXi-S2S forecast $c$. The fields are first compressed by an autoencoder~\cite{rombach2022ldm}, and a network $v_\theta$ learns
to transport noise $z_0$ to the observed field $z_1$ by minimising
\begin{equation}
\mathcal{L}(\theta)=\mathbb{E}_{t,z_0,z_1}\big\|\,v_\theta\big((1-t)z_0+t z_1,\,t,\,c\big)-(z_1-z_0)\big\|^2 .
\end{equation}
Sampling the trained model many times gives an ensemble with realistic spatial structure and a consistent rainfall--temperature relation.
A similar approach was recently used for S2S rainfall~\cite{shi2026generative}. Table~\ref{tab:tasks} lists the tasks.

\begin{table}[htb]
\centering\small
\caption{Tasks for the remainder of the thesis.}
\label{tab:tasks}
\renewcommand{\arraystretch}{1.12}
\begin{tabularx}{\linewidth}{@{}c >{\raggedright\arraybackslash}p{0.27\linewidth} >{\raggedright\arraybackslash}X@{}}
\toprule
& Task & Description and expected outcome \\
\midrule
T1 & Global probabilistic model & Complete the ECMWF AI Weather Quest model~\cite{aiwq2025}, which turns the FuXi-S2S ensemble into
five-category rainfall probabilities for days 19--25 and 26--32 and refines them with a U-Net/Transformer (about 150,000 parameters).
Early tests against IMERG give 3--4\% RPSS over climatology; the full ERA5 evaluation is still to be done. \\
T2 & Latent representation & Train an autoencoder that compresses weekly rainfall and temperature fields over India while keeping
heavy-rainfall structure. \\
T3 & Flow-matching model & Train the conditional flow-matching model in the latent space, first for rainfall and then jointly for
rainfall and temperature. \\
T4 & Extend the adapter & Apply the regional adapter to 2\,m temperature and to other forecast systems, to test whether it generalises
beyond one model and one variable. \\
T5 & Evaluation & Compare all new methods with bias correction, EMOS~\cite{gneiting2005emos} and the regional adapter, using CRPS,
reliability and coverage against both IMD and IMERG. The 2025 monsoon is kept aside as the final test. \\
T6 & Dissemination & Add calibrated forecasts to \url{s2s.scdlds.com}, submit the multi-model paper~\cite{raj2026piggy}, and write
the thesis. \\
\bottomrule
\end{tabularx}
\end{table}

% ============================ 5. SUMMARY ============================
\section{Summary}
In the first half of the project we have built a benchmark that shows where and when each forecast system is skilful over India, shown
that a small regional model can make the FuXi-S2S ensemble much more reliable, found that learned model combination is hard to separate
from simple bias correction, and put the forecasts into real-time use. The second half will develop a generative model for joint,
calibrated probabilistic forecasts of rainfall and temperature, evaluated against these baselines.

\section*{Acknowledgements}
I thank my supervisor Prof.\ Sandeep Kumar Juneja and my expert Dr.\ Bedartha Goswami for their guidance, and Dr.\ Manmeet Singh
(Western Kentucky University) and Dr.\ Parthasarathi Mukhopadhyay (Ashoka University) for their collaboration on this work.

% ============================ REFERENCES ============================
\clearpage
\renewcommand{\refname}{References}
\footnotesize
\setlength{\bibsep}{0.5pt}
\begin{multicols}{2}
\raggedright\sloppy
\begin{thebibliography}{99}
\bibitem{raj2026bench} A.~Raj, S.~Dhanuka, M.~Singh, P.~Mukhopadhyay, B.~Goswami, S.~K.~Juneja. IndiaS2S-Bench: A comprehensive benchmark of subseasonal to seasonal forecasting over India. \emph{Tackling Climate Change with ML workshop, NeurIPS}, 2026.
\bibitem{raj2026adapter} A.~Raj, M.~Singh, S.~Dhanuka, B.~Goswami, P.~Mukhopadhyay, S.~K.~Juneja, M.~Mohapatra. Calibrating AI subseasonal to seasonal (S2S) forecasts with lightweight regional adapters. \emph{Tackling Climate Change with ML workshop, NeurIPS}, 2026.
\bibitem{raj2026piggy} A.~Raj et al. When does lightweight post-processing help? A seven-system JJAS rainfall benchmark over India. Manuscript in preparation, 2026.
\bibitem{vitart2017s2s} F.~Vitart et al. The subseasonal to seasonal (S2S) prediction project database. \emph{Bull. Amer. Meteor. Soc.}, 98(1):163--173, 2017.
\bibitem{imd2022lpa} India Meteorological Department. Long range forecast for the 2022 southwest monsoon season rainfall (LPA 1971--2020 of 87~cm), 2022.
\bibitem{lam2023graphcast} R.~Lam et al. Learning skillful medium-range global weather forecasting. \emph{Science}, 382:1416--1421, 2023.
\bibitem{chen2024fuxis2s} L.~Chen et al. A machine learning model that outperforms conventional global subseasonal forecast models. \emph{Nat. Commun.}, 15:6425, 2024.
\bibitem{kochkov2024neuralgcm} D.~Kochkov et al. Neural general circulation models for weather and climate. \emph{Nature}, 632:1060--1066, 2024.
\bibitem{rasp2024wb2} S.~Rasp et al. WeatherBench 2: A benchmark for the next generation of data-driven global weather models. \emph{J. Adv. Model. Earth Syst.}, 16:e2023MS004019, 2024.
\bibitem{malik2024s2s} I.~Malik, V.~Mishra. Sub-seasonal to seasonal (S2S) prediction of dry and wet extremes for climate adaptation in India. \emph{Clim. Serv.}, 34:100457, 2024.
\bibitem{huffman2023imerg} G.~J.~Huffman et al. GPM IMERG Final precipitation L3 1 day 0.1$^\circ$ V07. NASA GES DISC, 2023.
\bibitem{pai2014imd} D.~S.~Pai et al. Development of a new high spatial resolution (0.25$^\circ$) long period (1901--2010) daily gridded rainfall data set over India. \emph{Mausam}, 65(1):1--18, 2014.
\bibitem{hersbach2020era5} H.~Hersbach et al. The ERA5 global reanalysis. \emph{Q. J. R. Meteorol. Soc.}, 146:1999--2049, 2020.
\bibitem{kimani2025piggycast} J.~Kimani, O.~Ang\'elil, C.~Toumping, S.~Knoblauch. PiggyCast: improving weather prediction accuracy through a stacking-based ensemble AI approach. \emph{Tackling Climate Change with ML workshop, NeurIPS}, 2025.
\bibitem{chen2016xgboost} T.~Chen, C.~Guestrin. XGBoost: A scalable tree boosting system. \emph{Proc. 22nd ACM SIGKDD}, 785--794, 2016.
\bibitem{aiwq2025} ECMWF. AI Weather Quest: a real-time sub-seasonal forecasting competition. \url{https://aiweatherquest.ecmwf.int}, 2025.
\bibitem{lipman2023flow} Y.~Lipman, R.~T.~Q.~Chen, H.~Ben-Hamu, M.~Nickel, M.~Le. Flow matching for generative modeling. \emph{ICLR}, 2023.
\bibitem{liu2023rectified} X.~Liu, C.~Gong, Q.~Liu. Flow straight and fast: learning to generate and transfer data with rectified flow. \emph{ICLR}, 2023.
\bibitem{rombach2022ldm} R.~Rombach et al. High-resolution image synthesis with latent diffusion models. \emph{CVPR}, 2022.
\bibitem{shi2026generative} W.~Shi et al. Distribution-guided ensemble postprocessing for S2S precipitation forecasts: a seamless pathway using deep generative models. \emph{JGR: Mach. Learn. Comput.}, 3:e2025JH000993, 2026.
\bibitem{gneiting2005emos} T.~Gneiting, A.~E.~Raftery, A.~H.~Westveld, T.~Goldman. Calibrated probabilistic forecasting using ensemble model output statistics and minimum CRPS estimation. \emph{Mon. Wea. Rev.}, 133:1098--1118, 2005.
\end{thebibliography}
\end{multicols}

\end{document}
